\documentclass{article}
\usepackage{graphicx} % Required for inserting images

\title{COmplex Var}
\author{Andrés López Gónzalez}
\date{September 2024}

\begin{document}

Resuelve formalmente el siguiente problema de variable compleja (en todos los casos demuestra tu respuesta).

1. Considera la funcion $u(x,y) = x^3 − 3xy^2$, ¿existe una funcion armonica v : R^2 → R, tal que sea conjugada armonica de u?
a) En general, ¿dada u : R^2 → R armonica, existe siempre su conjugado armonico v : R^2 → R?

b) ¿Y si D es un subconjunto de R2, y u : D ⊂ R^2 → R tambien se puede decir que siempre existe el conjugado armonico v : D ⊂ R^2 → R?

c) Sea u(x, y) = ln(\sqrt{x^2 + y^2}), ¿existe su conjugada armonica en R^2\{0}?

\end{document}

Resuelve formalmente el siguiente problema de variable compleja (en todos los casos demuestra tu respuesta).

2.1) Construye una funcion homografica (transformacion de Moebius) que lleve a la circunferencia unitaria en el eje real y a su interior en el semiplano superior. Nota para tí: Si se te dificulta ver esto, la circunferencia unitaria y su interior me refiero a los {(x,y) : x^2 + y^2 \leq 1}, al semiplano superior me refiero a los {(x,y) : y \geq 0}

2.2) ¿Son ciertas las siguientes propiedades? Dados a, b, c ∈ C
a) (ab)^c = a^cb^c
b) (a^b)^c = a^{bc}
c) a^ba^c = a^{b+c}
Da una demostracion o un contraejemplo en cada inciso.

Resuelve formalmente el siguiente problema de variable compleja (en todos los casos demuestra tu respuesta).

Sea g_1(z) la rama de logaritmo que considera argumentos de z en el intervalo (−π, π), a) define otra rama de logaritmo g_2(z) cuyo dominio sea C \ {x + iy : y ≤ 0, x = 0}, y tal que g_1 y g_2 coincidan en el semiplano superior. b) Define tambien otra rama g_3 con un dominio adecuado tal que coincida con g_2  ́unicamente en el segundo y tercer cuadrante. 

c) Explica cual es el fallo en la siguiente secuencia.
\frac{-1}{1} = \frac{1}{-1} \\
\frac{\sqrt{-1}}{\sqrt{1}} = \frac{\sqrt{1}}{\sqrt{-1}} \\
\frac{i}{1} = \frac{1}{i} \\
i^2 = 1 \\
-1 = 1